\documentclass[11pt]{article}

\usepackage{amsmath,amssymb,amsthm}
\usepackage{hyperref}
\usepackage{enumitem}
\usepackage{booktabs}

\newcommand{\lean}[1]{\texttt{#1}}
\newcommand{\file}[1]{\texttt{#1}}

\title{A 3D Interface for Neural Networks in Lean:\\
Deterministic Dynamics, Stochastic Kernels, and Energy-Based Learning}
\author{(anonymous)}
\date{}

\begin{document}
\maketitle

\begin{abstract}
We present a Lean~4 development that consolidates several semantic views of ``neural networks''
into a small set of compositional interfaces:
(i) deterministic dynamics, (ii) stochastic dynamics via Markov kernels, and (iii) energy models.
These are packaged as a lightweight \emph{3D} vocabulary and connected by bridges to Mathlib's
probability theory (kernels, invariance) and to Gibbs specifications for infinite lattice systems.
On top of this interface, we formalize Boltzmann Machine learning in a measure-theoretic style
(positive/negative phases as measures; learning signals as expectation differences) and prove
finite exponential-family identities, including a vector-parameter gradient formula
\(\nabla \log Z(\theta) = -\mathbb{E}_{\theta}[\mathrm{stat}]\).
We instantiate the theory for a concrete SymmetricBinary Hopfield Hamiltonian and package the
canonical ``data statistics minus model statistics'' update as a verified learning-rule adapter.
\end{abstract}

\section{Contributions}
The key contributions are:
\begin{itemize}[leftmargin=2em]
\item \textbf{A minimal 3D interface layer} (\file{NeuralNetwork/ThreeD/Core.lean}) providing
  \lean{EnergyModel}, \lean{DeterministicDynamics}, \lean{StochasticDynamics}, and
  \lean{IndexedStochasticDynamics}.
\item \textbf{Stationarity/invariance infrastructure} for these interfaces
  (\file{NeuralNetwork/ThreeD/Invariant.lean}), including composition over schedules
  (\lean{kernelOfList}) and a stationarity preservation lemma.
\item \textbf{Bridges to Mathlib kernels and Gibbs specifications}:
  deterministic maps to kernels
  (\file{NeuralNetwork/ThreeD/Bridges/DeterministicKernel.lean}),
  finite MCMC matrices to kernels
  (\file{NeuralNetwork/ThreeD/Bridges/MCMCFinite.lean}),
  and Gibbs specifications/potentials to indexed stochastic dynamics
  (\file{NeuralNetwork/ThreeD/Bridges/GibbsMeasure.lean},
   \file{NeuralNetwork/ThreeD/Bridges/GibbsPotential.lean}).
\item \textbf{A principled Boltzmann learning layer}
  (\file{NeuralNetwork/ThreeD/BoltzmannLearning/}): phases as measures, statistics as
  measurable functions, and a learning-rule schema (\lean{BM.LearningRule}).
\item \textbf{Finite and vector-parameter Gibbs identities in Lean}:
  \lean{FiniteGibbs.hasDerivAt_logZ} and \lean{VectorGibbs.hasGradientAt_logZ}, plus the
  single-sample negative log-likelihood gradient theorem
  \lean{VectorGibbs.hasGradientAt_negLogLik}.
\item \textbf{Concrete Hopfield/SymmetricBinary instantiation}:
  an explicit parameter Hilbert space and sufficient statistics such that the Hopfield Hamiltonian
  matches the abstract Gibbs energy
  (\lean{SymmetricBinaryInstantiation.hamiltonian_eq_vectorGibbs_energy}),
  and a packaged coordinate-wise BM update rule
  (\lean{SymmetricBinaryInstantiation.learningRule}).
\end{itemize}

\section{Repository entrypoints and build}
Stable import surfaces are provided by:
\begin{itemize}[leftmargin=2em]
\item \file{NeuralNetwork/SOTA.lean} (project-level entrypoint),
\item \file{NeuralNetwork/ThreeD/SOTA.lean} and \file{NeuralNetwork/ThreeD/SOTA_Gibbs.lean},
\item \file{NeuralNetwork/ThreeD/BoltzmannLearning/SOTA.lean}.
\end{itemize}

Artifact build commands:
\begin{itemize}[leftmargin=2em]
\item \lean{lake build NeuralNetwork.SOTA}
\item \lean{lake build NeuralNetwork.ThreeD.BoltzmannLearning.SOTA}
\item \lean{lake build NeuralNetwork.ThreeD.SOTA_Gibbs}
\end{itemize}

\section{Main theorems (Lean names)}
\subsection{Finite scalar parameter}
\begin{itemize}[leftmargin=2em]
\item \lean{NeuralNetwork.ThreeD.BoltzmannLearning.FiniteGibbs.hasDerivAt_logZ}:
  for a finite type \(X\) and energy feature \(\varphi : X \to \mathbb{R}\),
  \(\frac{d}{dt} \log Z(t) = -\mathbb{E}_t[\varphi]\).
\end{itemize}

\subsection{Vector parameter in a Hilbert space}
\begin{itemize}[leftmargin=2em]
\item \lean{NeuralNetwork.ThreeD.BoltzmannLearning.VectorGibbs.hasGradientAt_logZ}:
  for \(\theta : \Theta\) in a real Hilbert space and statistics
  \(\mathrm{stat} : X \to \Theta\),
  \(\nabla \log Z(\theta) = -\mathbb{E}_{\theta}[\mathrm{stat}]\).
\item \lean{NeuralNetwork.ThreeD.BoltzmannLearning.VectorGibbs.hasGradientAt_negLogLik}:
  \(\nabla\mathrm{NLL}(\theta;x)=\mathrm{stat}(x)-\mathbb{E}_\theta[\mathrm{stat}]\).
\end{itemize}

\subsection{Concrete Hopfield/SymmetricBinary connection}
\begin{itemize}[leftmargin=2em]
\item \lean{NeuralNetwork.ThreeD.BoltzmannLearning.SymmetricBinaryInstantiation.hamiltonian_eq_vectorGibbs_energy}
  (\file{NeuralNetwork/ThreeD/BoltzmannLearning/SymmetricBinaryInstantiation.lean}).
\item \lean{NeuralNetwork.ThreeD.BoltzmannLearning.SymmetricBinaryInstantiation.learningRule}
  (\file{NeuralNetwork/ThreeD/BoltzmannLearning/SymmetricBinaryLearningRule.lean}).
\end{itemize}

\section{Scope and current limitations}
\begin{itemize}[leftmargin=2em]
\item \textbf{Finiteness vs infinite systems.} The Boltzmann learning theorems are currently
  proved in a finite state setting (finite sums) and for Hilbert-space parameters
  (Mathlib's gradient API). Separately, infinite-lattice semantics are represented via
  Gibbs specifications as indexed kernels; extending learning theorems to that setting requires
  additional analytic hypotheses (e.g.\ integrability/differentiation under the integral).
\item \textbf{Interface validation.} The 3D layer is intentionally small; its adequacy is currently
  demonstrated by multiple independent bridges (deterministic kernels, MCMC kernels, Gibbs specs,
  and energy-based learning instantiations).
\end{itemize}

\section{Module map (high-level)}
\begin{center}
\begin{tabular}{@{}ll@{}}
\toprule
Area & Key files \\
\midrule
3D core & \file{NeuralNetwork/ThreeD/Core.lean} \\
Stationarity & \file{NeuralNetwork/ThreeD/Invariant.lean} \\
Deterministic $\to$ kernel & \file{NeuralNetwork/ThreeD/Bridges/DeterministicKernel.lean} \\
Finite MCMC $\to$ kernel & \file{NeuralNetwork/ThreeD/Bridges/MCMCFinite.lean} \\
Gibbs specs/potentials & \file{NeuralNetwork/ThreeD/Bridges/GibbsMeasure.lean}, \file{.../GibbsPotential.lean} \\
BM learning core & \file{NeuralNetwork/ThreeD/BoltzmannLearning/Core.lean} \\
Finite Gibbs identity & \file{NeuralNetwork/ThreeD/BoltzmannLearning/FiniteGibbs.lean} \\
Vector Gibbs identity & \file{NeuralNetwork/ThreeD/BoltzmannLearning/VectorGibbs.lean} \\
NLL gradient & \file{NeuralNetwork/ThreeD/BoltzmannLearning/VectorGibbsLearning.lean} \\
Hopfield instantiation & \file{NeuralNetwork/ThreeD/BoltzmannLearning/SymmetricBinaryInstantiation.lean} \\
LearningRule adapter & \file{NeuralNetwork/ThreeD/BoltzmannLearning/SymmetricBinaryLearningRule.lean} \\
\bottomrule
\end{tabular}
\end{center}

\section*{Acknowledgements and bibliography}
See \file{refs.bib}. (Add the venue template's bibliography style as needed.)

\bibliographystyle{plain}
\bibliography{refs}

\end{document}

